\subsection{Solucion desarrollada}

\begin{respuesta}{Situacion actual}
La duracion de los bloques es \(6,2,7,9\) horas. Por lo tanto:
\[
\lambda=(5.83,\ 13.50,\ 12.14,\ 5.00)\ \text{viajes/hora}.
\]
Con \(\mu=3\):
\[
\rho_b=\frac{\lambda_b}{c_b\mu}.
\]
Los resultados son:
\begin{center}
\begin{tabular}{ccccc}
\toprule
Bloque & \(\lambda\) & Taxis & \(\rho\) & \(W_q\) min \\
\midrule
1 & 5.83 & 4 & 0.486 & 29.66 \\
2 & 13.50 & 6 & 0.750 & 87.86 \\
3 & 12.14 & 6 & 0.675 & 92.26 \\
4 & 5.00 & 3 & 0.556 & 8.32 \\
\bottomrule
\end{tabular}
\end{center}
El bloque con mayor espera es el 3; el de mayor utilizacion es el 2. La politica de no superar \(90\%\) se cumple en todos los bloques.
\end{respuesta}

\begin{respuesta}{Viajes fuera de zona}
La demanda total base es:
\[
35+27+85+45=192.
\]
El aumento es:
\[
\Delta=0.2(192)=38.4.
\]
Las nuevas demandas son:
\[
D'=(40.76,\ 32.76,\ 104.20,\ 52.68).
\]
Dividiendo por la duracion de cada bloque:
\[
\lambda'=(6.79,\ 16.38,\ 14.89,\ 5.85).
\]
Si fuera de zona reduce la tasa a \(\mu=1.5\) viajes/hora por taxi y se mantiene la flota original:
\[
\rho'=(1.132,\ 1.820,\ 1.654,\ 1.301).
\]
Todos los bloques quedan inestables. Para mantener \(\rho\le 0.9\):
\[
c_b\ge \left\lceil\frac{\lambda'_b}{0.9(1.5)}\right\rceil.
\]
Asi:
\[
c=(6,\ 13,\ 12,\ 5).
\]
La decision de salir fuera de zona requiere mucha mas flota; no debe aceptarse sin comparar margen adicional contra costo de capacidad y espera.
\end{respuesta}

\begin{respuesta}{Modelo matematico}
Variables:
\[
k_b\in\mathbb{Z}_+,\qquad z\in\{0,1\},\qquad \rho_b\ge 0,\qquad W_{q,b}\ge 0.
\]
La variable \(z=1\) indica aceptar viajes fuera de zona. La demanda por bloque es:
\[
\lambda_b(z)=\frac{D_b+z\Delta s_b}{H_b}.
\]
La utilizacion:
\[
\rho_b=\frac{\lambda_b(z)}{k_b\mu(z)}.
\]
Restricciones:
\[
\rho_b\le 0.9,\qquad k_b\in\mathbb{Z}_+,\qquad z\in\{0,1\}.
\]
Funcion objetivo:
\[
\min \sum_b C_q\lambda_b(z)W_{q,b}
+\sum_b C_u(0.9-\rho_b)
+\sum_b C_k k_b.
\]
La espera se aproxima imponiendo:
\[
W_{q,b}\ge
\left(\frac{CV_{a,b}^2+CV_{s,b}^2}{2}\right)
\frac{\rho_b^{\sqrt{2k_b+2}-1}}{k_b(1-\rho_b)\mu(z)}.
\]
El modelo decide simultaneamente si se aceptan viajes fuera de zona y cuantos taxis asignar a cada bloque.
\end{respuesta}
